\section{Alphateilchen im Feld der Spulen: Bleiben sie lange genug im Plasma?}\label{sec:alpha}

Im Abschnitt über das Plasma (Abschnitt~\ref{sec:plasma}) heizen die Alphateilchen das Plasma mit \SI{611}{MW}. Das
setzt voraus, dass jedes Alpha seine Energie im Plasma abgibt, bevor es den Rand erreicht. Ob das stimmt, hängt vom
Magnetfeld ab. Der Nutzer fragte nach einer dynamischen Simulation mit dem Magnetfeld; dies ist der erste Teil davon
(Variante A): die Bahnen von 1000 Alphas im Feld der eigenen Spulen.

\paragraph{Begriffe in diesem Abschnitt.}
\begin{description}\itemsep2pt
\item[Alphateilchen] der Heliumkern, der bei jeder Fusion von Deuterium und Tritium entsteht, mit \SI{3,5}{MeV}
  Bewegungsenergie. Er ist elektrisch geladen und folgt deshalb den Magnetfeldlinien.
\item[Führungszentrum] Ein geladenes Teilchen kreist schnell um eine Feldlinie (beim Alpha mit etwa \SI{4}{cm}
  Radius). Statt jeden Kreis zu rechnen, verfolgt man nur den Mittelpunkt dieses Kreises, das Führungszentrum. Es
  läuft an der Feldlinie entlang und driftet langsam quer dazu.
\item[umlaufend und gefangen] Ein Alpha, das überwiegend entlang des Feldes fliegt, läuft immer wieder um den ganzen
  Ring (\emph{umlaufend}). Eines, das überwiegend quer zum Feld fliegt, wird an Stellen mit stärkerem Feld
  zurückgeworfen wie von einem Spiegel und pendelt hin und her (\emph{gefangen}). Gefangene Teilchen sind besonders
  empfindlich für kleine Fehler im Feld.
\item[$v_\parallel/v$] der Anteil der Geschwindigkeit entlang des Feldes: $+1$ oder $-1$ heißt genau entlang,
  $0$ genau quer.
\item[Abbremszeit] die Zeit, in der ein Alpha durch Stöße mit den Elektronen und Ionen seine Energie an das Plasma
  abgibt, hier etwa \SI{0,2}{s}. Was so lange im Plasma bleibt, hat seine Aufgabe erfüllt.
\item[letzte geschlossene Fläche] die äußerste der ineinanderliegenden Flussflächen, der Rand des eingeschlossenen
  Plasmas. Wer sie überquert, gilt hier als verloren.
\item[Quasisymmetrie] eine besondere Form des Stellaratorfelds: Obwohl die Anlage verdreht ist, sieht die Feldstärke
  entlang der Bahnen so aus wie in einem symmetrischen Ring. Dann driften gefangene Teilchen im Mittel nicht nach
  außen. Unsere Plasmaform ist quasi-axialsymmetrisch (QA) ausgelegt.
\item[Prompter Verlust] ein Alpha, dessen Bahn schon beim ersten Umlauf den Rand streift; es geht innerhalb von
  Bruchteilen einer Millisekunde verloren. Das betrifft vor allem Alphas, die nahe am Rand geboren werden.
\end{description}

\paragraph{Methode.} Die Bahnen werden mit simsopt \cite{simsopt} als stoßfreie Führungszentren im Vakuumfeld der
16 Spulen aus Stufe~2 gerechnet. Spulen und Form liegen in Einheitsgröße vor ($R_0 = \SI{1,03}{m}$). Stoßfreie Bahnen
hängen aber nur vom Verhältnis Kreisradius zu Plasmaradius ab ($\rho_L/a = 0{,}018$). Statt die Spulen zu vergrößern,
wird deshalb die Geschwindigkeit passend verkleinert, und Zeiten und Orte werden auf die Anlage zurückgerechnet. Die
1000 Startorte folgen der Verteilung der Fusionsleistung, die Richtungen sind gleichmäßig verteilt. Gerechnet wird bis
\SI{0,2}{s}, der Abbremszeit. Für die Energiebilanz bekommt jedes verlorene Alpha die Energie, die es zum
Verlustzeitpunkt nach der klassischen Abbremsung \cite{stix1972} noch hätte:
\[
u(t) = \bigl(u_0 + 1\bigr)\,e^{-3t/\tau_s} - 1,\qquad u = (E/E_c)^{3/2},\qquad \tau_s = \SI{0,195}{s},\ E_c = \SI{415}{keV}.
\]
Die Verlustorte werden einem Gitter auf der Randfläche zugeordnet ($64\times32$ Punkte, zusammengefasst zu
$16\times8$ Feldern von je etwa \SI{10}{m^2}). Daraus folgt der Wärmefluss.

\bild{alpha_verluste.png}{1.0}{Verluste der Alphateilchen: Zeitverlauf, wer verloren geht, wo sie den Rand treffen,
Kennzahlen (1000 Alphas, stoßfrei, Feld der eigenen Spulen).}
\lesehilfe{Oben links: Waagrecht die Zeit seit der Geburt der Alphas, logarithmisch von einer Hundertstel
Millisekunde bis \SI{0,2}{s}. Senkrecht, wie viel Prozent verloren sind. Blau zählt Teilchen, orange ihre Energie.
Die orange Kurve bleibt darunter, weil spät verlorene Alphas schon einen großen Teil ihrer Energie abgegeben haben.
In der ersten Millisekunde geht kaum etwas verloren; die meisten Verluste kommen zwischen \SI{5}{ms} und
\SI{100}{ms}. Oben rechts: Jeder Punkt ist ein Alpha, waagrecht seine Flugrichtung bei der Geburt (0 = quer zum Feld),
senkrecht sein Geburtsort (0 = Plasmamitte, 1 = Rand). Die roten, verlorenen Alphas liegen fast alle in einem
schmalen Streifen um 0: Es sind gefangene Teilchen. Unten links: die ausgerollte Randfläche des Plasmas, waagrecht
einmal den Ring entlang, senkrecht einmal um den Querschnitt. Die Farbe ist der Wärmefluss, die Zahl die Anzahl der
Treffer. Die Treffer liegen gebündelt in einem Band, und zwar zweimal, weil sich die Anlage zweimal wiederholt.
Unten rechts die wichtigsten Zahlen.}

\bild{alpha_bahnen_standbild.png}{1.0}{Drei Alphabahnen im Feld der eigenen Spulen (Standbild aus dem Video
\texttt{docs/filme/alpha\_bahnen.mp4}): umlaufend (blau), gefangen (gelb), früh verloren (rot).}
\lesehilfe{Man blickt schräg von oben auf den Ring. Die dünnen grauen Schleifen sind Querschnitte des Plasmarands an
16 Stellen; sie zeigen, wie sich die Form beim Umlauf verdreht. Jede farbige Linie ist der Weg des Führungszentrums
eines Alphas. Das blaue Alpha fliegt entlang des Feldes und läuft immer wieder um den ganzen Ring; es bleibt sicher
innen. Das gelbe fliegt fast quer zum Feld und kehrt an Stellen mit starkem Feld um; seine Bahn ist an den Enden
umgebogen. Das rote ist ebenfalls ein gefangenes Alpha. Seine Bahn wandert langsam nach außen, bis es nach
\SI{0,13}{ms} den Rand erreicht (roter Punkt oben). Solche Bahnen erzeugen die Verluste im Bild oben, meist allerdings viel langsamer.}

\paragraph{Ergebnisse.}
\begin{itemize}\itemsep1pt
\item \textbf{16,6\,\% der Alphas gehen verloren} ($\pm1{,}2$\,\% statistische Unsicherheit bei 1000 Teilchen),
  davon nur 10 in der ersten Millisekunde. Die Hälfte der Verluste ist nach \SI{25}{ms} erreicht.
\item \textbf{Es sind fast nur gefangene Alphas}: 95\,\% der verlorenen hatten $|v_\parallel/v| < 0{,}3$. Sie stammen
  im Mittel etwas weiter außen ($s = 0{,}24$ gegen $0{,}18$).
\item \textbf{Verlorene Energie: 11,6\,\% der Alphaleistung, also \SI{71}{MW}}. Dem Plasma fehlt dieser Teil seiner
  Selbstheizung; statt \SI{611}{MW} bleiben etwa \SI{540}{MW}.
\item \textbf{Die Verluste treffen gebündelt auf.} Im Mittel über den ganzen Rand sind es nur \SI{58}{kW/m^2}, aber
  fast alle Treffer liegen in einem Band in der oberen Hälfte des Querschnitts, in zwei Flecken. Im heißesten Feld
  von etwa \SI{10}{m^2} sind es rund \SI{2,4}{MW/m^2}. Die Erstwand dieses Projekts ist auf \SI{0,5}{MW/m^2}
  ausgelegt (Abschnitt Kühlung). Wegen der groben Felder und der nur 27 Treffer im heißesten Feld ist der Spitzenwert
  unsicher; die echte Spitze ist eher höher.
\end{itemize}

\paragraph{Einordnung.} Die Verluste kommen nicht als schneller Verlust in den ersten Umläufen, sondern langsam, über
Millisekunden bis zur Abbremszeit. Das ist das typische Bild, wenn gefangene Teilchen in einem nicht ganz
quasisymmetrischen Feld nach außen driften. Die Plasmaform selbst ist sehr genau quasisymmetrisch (Zielwert
der Quasisymmetrie in DESC $2{,}4\cdot10^{-5}$, kleiner als bei der Referenz \emph{precise QA} mit
$5{,}8\cdot10^{-5}$). Die Spulen treffen diese Form aber nur ungefähr: Das Feld steht im Mittel um 0,27\,\% und
höchstens um 0,96\,\% schräg zur Randfläche ($|B_n|/|B|$). Ob der Hauptteil der Verluste aus diesen Spulenfehlern stammt oder schon in der Form angelegt ist,
klärt der folgende Vergleich. Ein Verlust von über 10\,\% der Alphaenergie wäre für ein Kraftwerk zu viel, vor allem
wegen der gebündelten Wandlast.

\paragraph{Vorher und nachher: Spulen oder Form?} Dieselbe Rechnung wurde ein zweites Mal gemacht, diesmal im Feld der
Plasmaform allein, also so, als ob die Spulen die Form perfekt träfen. Dazu wird das Gleichgewicht der Form mit
\texttt{booz\_xform}\footnote{\url{https://github.com/hiddenSymmetries/booz_xform}} in
\emph{Boozer-Koordinaten} umgerechnet, ein Koordinatensystem, in dem die Feldlinien gerade verlaufen und die
Feldstärke als einfache Fourier-Reihe vorliegt. Dann rechnet simsopt die Bahnen direkt darin, mit derselben
Verteilung der Geburtsorte und Richtungen wie zuvor (im Mittel $s = 0{,}18$, 31\,\% gefangene Startrichtungen in
beiden Läufen).

\emph{Erwartung vorher.} Bevor das Ergebnis vorlag, wurde im Ereignisprotokoll festgelegt: Die Form ist genauer
quasisymmetrisch als \emph{precise QA}, für die Landreman und Paul in \SI{0,2}{s} keine Alphaverluste fanden
\cite{landreman2022}. Erwartet wurden deshalb unter 1\,\% Verlust. Die Deutung wurde ebenfalls vorab festgelegt:
unter 2\,\% heißt, die Spulen sind die Hauptursache; zwischen 2 und 10\,\% tragen beide bei; über 10\,\% liegt es an
der Form selbst oder an einem Rechenfehler.

\bild{alpha_vergleich.png}{1.0}{Vorher und nachher: Alphaverluste im Feld der 16 Spulen (rot) und im Feld der
Plasmaform allein (blau), je 1000 Alphas bis \SI{0,2}{s}.}
\lesehilfe{Links: waagrecht die Zeit seit der Geburt der Alphas (logarithmisch), senkrecht der verlorene Anteil. Die
rote Kurve (Spulen) steigt ab etwa einer Millisekunde stetig bis 16,6\,\%. Die blaue Kurve (Form allein) springt ganz
am Anfang auf 0,3\,\% und bleibt dann flach: Danach geht kein einziges Alpha mehr verloren. Die gestrichelte Linie
ist die vorher festgelegte Erwartung von 1\,\%; die blaue Kurve bleibt darunter. Rechts die Kennzahlen beider Läufe.}

\emph{Ergebnis.} Im Feld der Form allein gehen nur 3 von 1000 Alphas verloren ($0{,}3 \pm 0{,}2$\,\%), alle weit
außen geboren ($s = 0{,}60$ bis $0{,}72$) und schon in den ersten \SI{0,1}{ms}. Die langsamen Verluste gefangener
Alphas, die im Spulenfeld 95\,\% der Verluste ausmachen, fehlen vollständig. Die verlorene Heizleistung sinkt von
\SI{71}{MW} auf \SI{1,8}{MW}. Weil dieser Lauf viel schneller war als der Spulenlauf, wurde stichprobenhaft
geprüft, ob die Bahnen wirklich bis \SI{0,2}{s} gerechnet werden: Sechs nicht verlorene Alphas (vier gefangene, zwei
umlaufende) erreichen alle \SI{0,2000}{s}; die gefangenen pendeln (v$_\parallel$ wechselt immer wieder das
Vorzeichen) und bleiben innen. Die Erwartung ist damit bestätigt: Rund 98\,\% der Verluste stammen von der
Ungenauigkeit der Spulen, nicht von der Plasmaform. Das passt zum Standbild oben: Die rote Bahn bleibt in einem
kleinen Feldtal hängen, wie es nur ein nicht ganz quasisymmetrisches Feld bildet. Dieser Vergleich ist die
Ausgangslage für den nächsten Schritt: Spulen nachoptimieren und dieselbe Rechnung wiederholen.

\paragraph{Grenzen.}
\begin{itemize}\itemsep1pt
\item \textbf{Stoßfrei}: Die Bahnen ändern sich in Wirklichkeit, während das Alpha abbremst und durch Stöße die
  Richtung wechselt. Stöße können gefangene Teilchen in umlaufende verwandeln und umgekehrt; das kann die Verluste
  in beide Richtungen verschieben.
\item \textbf{Vakuumfeld}: Ströme im Plasma (endlicher Druck) sind nicht enthalten.
\item \textbf{Verlustort ist der Plasmarand}, nicht die Wand. Die Wand liegt weiter außen; die Treffer verteilen sich
  dort etwas anders.
\item \textbf{Statistik}: 1000 Alphas reichen für den Verlustanteil, aber nicht für eine feine Wandlastkarte.
\item \textbf{Feld der Form}: Die Umrechnung in Boozer-Koordinaten nutzt 24 Fourier-Moden in jeder Richtung; feinere
  Anteile fehlen. Das könnte die Verluste der Form leicht unterschätzen, ändert aber nicht die Größenordnung
  (0,3\,\% gegen 16,6\,\%).
\end{itemize}

\paragraph{Nächste Schritte.} (1) Erledigt: Die Ursache ist getrennt, es sind die Spulen. (2) Die Spulen genauer an
die Form anpassen, mit einem Ziel für Alphaverluste, und denselben Vergleich wiederholen („nachher“).
(3) Mehr Teilchen und Stöße für eine belastbare Wandlastkarte.
